\chapter{Coloboma}
\index{Coloboma}
\label{sec:Coloboma}

\section{Overview}
\begin{itemize}[noitemsep]
  \item Coloboma is a congenital eye malformation in which part of the eye tissue is missing, caused by incomplete closure of the embryonic eye (optic fissure) during development
  \item It most often affects the iris ("keyhole pupil"), but can involve the retina, optic nerve, lens, or eyelid
  \item It affects roughly 1 in 10,000 births and may occur in one or both eyes
\end{itemize}
\section{Causes}
This condition happens because of failure of the optic fissure to close fully during early fetal development (around the 5th--7th week of gestation). It can be inherited in autosomal dominant, recessive, or X-linked patterns, or arise from new mutations. It is also associated with chromosomal abnormalities and syndromes such as CHARGE syndrome and cat-eye syndrome.
\section{Risk Factors}
\begin{itemize}[noitemsep]
  \item Family history of coloboma
  \item Genetic syndromes (CHARGE syndrome, cat-eye syndrome)
  \item Chromosomal abnormalities
  \item Certain genetic mutations (e.g., PAX2, PAX6, CHD7)
  \item Exposure to some medications or toxins in early pregnancy
\end{itemize}
\section{Signs and Symptoms}
\subsection{Common symptoms}
\begin{itemize}[noitemsep]
  \item Iris coloboma: a keyhole or cat-eye shaped pupil
  \item Retinal or optic nerve coloboma: reduced vision, visual field defects
  \item Light sensitivity (photophobia) with iris involvement
  \item May be associated with a notch in the eyelid (eyelid coloboma)
  \item Some people have no visual symptoms if the defect is peripheral
\end{itemize}
\subsection{Emergency warning signs}
\begin{itemize}[noitemsep]
  \item Sudden visual change, flashes, floaters, or a shadow in the vision requires urgent evaluation, as coloboma increases the risk of retinal detachment
\end{itemize}
\section{Progression and Stages}
Coloboma is a structural condition present from birth and is generally stable. Visual function depends on which part of the eye is affected and its size. Complications such as retinal detachment, glaucoma, or cataract can develop later and require monitoring. Associated systemic anomalies in syndromic cases may dominate the clinical picture.
\section{Complications}
\begin{itemize}[noitemsep]
  \item Retinal detachment (increased risk)
  \item Glaucoma
  \item Cataract formation
  \item Strabismus (misaligned eyes) and amblyopia (lazy eye)
  \item Reduced vision or blindness in the affected eye
  \item Associated systemic abnormalities in syndromic forms
\end{itemize}
\section{Diagnosis}
\begin{itemize}[noitemsep]
  \item Eye examination confirming the characteristic missing tissue
  \item Slit-lamp examination for iris, lens, and anterior chamber involvement
  \item Dilated fundus examination for retinal and optic nerve coloboma
  \item Optical coherence tomography (OCT) and fundus photography
  \item Genetic testing and evaluation for associated syndromes
  \item Ultrasound or imaging of the eye if the view is obstructed
\end{itemize}
\section{Prevention}
\begin{itemize}[noitemsep]
  \item Coloboma is congenital and largely not preventable
  \item Genetic counseling for affected families
  \item Prenatal diagnosis in known familial cases
  \item Early eye examination after birth to detect and manage complications
  \item Avoidance of known teratogenic exposures in pregnancy
\end{itemize}
\section{Treatment Options}
\begin{itemize}[noitemsep]
  \item Cosmetic contact lenses or surgery for iris coloboma
  \item Correction of refractive error and treatment of amblyopia and strabismus
  \item Monitoring and treatment of glaucoma and cataract
  \item Timely treatment of retinal detachment, which is more challenging in coloboma
  \item Low-vision support for significant visual impairment
  \item Management of associated systemic conditions in syndromic cases
\end{itemize}
\section{Living with It}
\begin{itemize}[noitemsep]
  \item Attend regular eye examinations to monitor for complications
  \item Protect the eyes and report sudden visual changes promptly
  \item Use sunglasses for light sensitivity
  \item Seek support for cosmetic or functional concerns
  \item Coordinate care with a multidisciplinary team in syndromic cases
\end{itemize}
\section{Outlook}
Coloboma is stable, and most people retain good vision when the macula and optic nerve are not involved. With monitoring and treatment of complications, the long-term outlook is generally good, though retinal detachment remains a lifelong risk that requires vigilance.
